% \section{Success Story} \label{sec:deploymentinindustry}
\section{An Industrial Experience Report} \label{sec:deploymentinindustry}

%\subsection{Industrial system}
System X is a commercial software system that provides government-regulation related reporting services. The service is widely used as the market leader of the domain. 
%Due to a Non-Disclosure Agreement (NDA), we cannot reveal additional details about the system. We do note that 
System X has over ten years of history with more than two million lines of code that are based on Microsoft .Net. 
System X is deployed in an internal production environment and is used by external enterprise customers worldwide. 
%In particular, CPU usage is the main concern in performance regression detection.
%We retrieved the execution logs and the corresponding CPU usage on a daily basis. 
%System X is deployed in an internal production environment. 
System X is a CPU-intensive system. 
Thus, CPU usage is the main concern of our industrial collaborator in performance regression detection.
%We retrieved the execution logs and the CPU usage on a daily basis. 
Due to a Non-Disclosure Agreement (NDA), we cannot reveal additional details about the hardware environment and the usage scenarios of System X.

% \subsection{Industrial motivation}
% \lizhi{ref. R1.2 Our approach is motivated by the real problem encountered by our industrial collaborator. Our industrial collaborator uses an agile development approach and often releases a new version of their system within one week. Hence, they do not have enough budget and time to conduct performance tests. Therefore, the industrial collaborator looks for technologies that can tell them whether the system has a performance regression after the system has been pushed into the field and under real-user workloads, while such workloads are not identical from version to version. In our approach, we aim at directly detecting performance regressions based on the field operations of software systems. Such automate detection can complement or even replace typical in-house performance testing when testing resources are limited (e.g., in an agile environment) and/or field workloads are challenging to be reproduced in the testing environment. }

\subsection{Modeling system performance}

%For System X, 
We directly use the field data that is generated by workloads from the real end users to model the performance of System X. For each release cycle with $n$ days, we use the data from the first $n/2$ days to build the models and apply them on the data from the second $n/2$ days to evaluate the prediction performance of the models. We would like to note that there exists no control on the end users for applying any particular workload, and that all the data is directly retrieved from the field with no interference on the behavior of the end users. We studied 10 releases of System X.

%Due to an NDA, the detailed results of using the machine learning techniques to model the performance of the industry system (i.e., System X) are not presented in this paper. However, the results are similar to the shown results for the open source systems. In particular,
%When modeling the performance of all 10 studied versions of System X, the prediction errors are all between 10\% to 20\%. In addition, when we evaluate the models with different workloads, the differences between prediction errors are all either statistically insignificant or with negligible/small effect sizes.
We find that black-box models can accurately capture the relationship between the system performance and the system activities recorded in the logs, with a prediction error between 10\% to 20\%.
In particular, the models trained from the first half-release data can effectively predict the system performance of the second half release.

\subsection{Detecting performance regressions}

%For System X from the industry, 
For every new release, we use our approaches to compare the field data from the new release and the previous release to determine whether there are performance regressions. Since there are no injected or pre-known performance regressions, we present the detected performance regressions to the developers of the system and manually study the code changes to understand whether the detection results are correct. For the releases that are detected as not having performance regressions, we cannot guarantee that these releases are free of performance regressions. However, we also present the results to the developers to confirm whether there exist any users who report performance-related issues for these releases. If so, our detection results may be considered false.

Within all the 10 studied releases of System X, our approaches detected performance regressions from one release.
All three approaches detected performance regressions from the release with large effect sizes when compared with the previous release.
None of the three approaches false-positively detect performance regressions from the other nine releases (i.e., with either statistically insignificant difference or negligible effect sizes when compared with the previous release). 
By further investigating the release with detected performance regressions, we observed that developers added a synchronized operation to lock the resources that are responsible for generating a report, in order to protect the shared resources under the multi-thread situation. 
However, the reporting process is rather resource-heavy, resulting in significant overhead for each thread to wait and acquire the resources. Thus, the newly added lock causes the performance regression. After we discussed with the developers who are responsible for this module, we confirmed that this synchronized operation introduced the performance regression to the software. 
In addition, for all the nine releases from which our approaches did not detect performance regressions, the developers of System X have not yet received any reported performance issues from the end users till the day of writing this paper.

Our approaches have been adopted by our industrial collaborator to detect performance regressions of System X on a daily basis. 